\documentclass[10pt,a4paper]{article}

\usepackage[a4paper,top=14mm,bottom=16mm,left=14mm,right=14mm]{geometry}
\usepackage{fontspec}
\usepackage[table]{xcolor}
\usepackage{tabularx}
\usepackage{array}
\usepackage{enumitem}
\usepackage{titlesec}
\usepackage{fancyhdr}
\usepackage{lastpage}
\usepackage{ragged2e}
\usepackage{amssymb}
\usepackage[most]{tcolorbox}
\usepackage[hidelinks]{hyperref}

\IfFontExistsTF{Arial}{
  \setmainfont{Arial}
  \setsansfont{Arial}
}{
  \setmainfont{TeX Gyre Heros}
  \setsansfont{TeX Gyre Heros}
}

\definecolor{AIEPRed}{HTML}{D62027}
\definecolor{AIEPNavy}{HTML}{082B5C}
\definecolor{AIEPInk}{HTML}{243B53}
\definecolor{AIEPSlate}{HTML}{52606D}
\definecolor{AIEPBlue}{HTML}{1B6CA8}
\definecolor{AIEPCyan}{HTML}{35B7C6}
\definecolor{AIEPGreen}{HTML}{2E8B57}
\definecolor{AIEPGold}{HTML}{D8A928}
\definecolor{AIEPSoftBlue}{HTML}{E8F1F8}
\definecolor{AIEPGreenSoft}{HTML}{EAF5EF}
\definecolor{AIEPGoldSoft}{HTML}{FFF5D8}
\definecolor{AIEPRedSoft}{HTML}{FCEBED}
\definecolor{Line}{HTML}{D2DAE3}
\definecolor{White}{HTML}{FFFFFF}

\pagestyle{fancy}
\fancyhf{}
\renewcommand{\headrulewidth}{0pt}
\fancyfoot[L]{\scriptsize\textcolor{AIEPSlate}{Concurso Final GeoGreen Escolar 2026}}
\fancyfoot[R]{\scriptsize\textcolor{AIEPSlate}{AIEP Osorno · \thepage/\pageref{LastPage}}}

\setlength{\parindent}{0pt}
\setlength{\parskip}{3.5pt}
\setlist[itemize]{leftmargin=4.4mm,topsep=1pt,itemsep=1.1pt,parsep=0pt}
\setlist[enumerate]{leftmargin=5.2mm,topsep=1pt,itemsep=1.2pt,parsep=0pt}
\renewcommand{\arraystretch}{1.18}
\titleformat{\section}{\Large\bfseries\color{AIEPNavy}}{}{0pt}{}
\titleformat{\subsection}{\normalsize\bfseries\color{AIEPInk}}{}{0pt}{}

\newcolumntype{Y}{>{\RaggedRight\arraybackslash}X}
\newcolumntype{L}[1]{>{\RaggedRight\arraybackslash}p{#1}}
\newcolumntype{C}[1]{>{\centering\arraybackslash}p{#1}}

\newtcolorbox{card}[1][]{
  enhanced,colback=White,colframe=Line,boxrule=.5pt,arc=1.2mm,
  left=3.2mm,right=3.2mm,top=2.7mm,bottom=2.7mm,#1
}
\newtcolorbox{softcard}[2][]{
  enhanced,colback=#2!10,colframe=#2,boxrule=.55pt,arc=1.2mm,
  left=3.2mm,right=3.2mm,top=2.7mm,bottom=2.7mm,#1
}
\newcommand{\kicker}[1]{%
  {\color{AIEPRed}\scriptsize\bfseries\MakeUppercase{#1}}\par\vspace{2.4mm}
  {\color{AIEPRed}\rule{13mm}{1.2pt}}\vspace{3mm}
}
\newcommand{\labeltext}[1]{%
  {\scriptsize\bfseries\color{AIEPSlate}\MakeUppercase{#1}}\par\vspace{1.2mm}
}
\newcommand{\level}[2]{%
  \fcolorbox{#1}{#1}{\textcolor{White}{\bfseries\strut #2}}%
}

\begin{document}

% 1 ---------------------------------------------------------------------------
\thispagestyle{fancy}
\kicker{GeoGreen Escolar · AIEP Osorno}
{\Huge\bfseries\color{AIEPNavy} Concurso Final GeoGreen Escolar 2026}\par
\vspace{2mm}
{\large\color{AIEPSlate} Lineamientos de participación, presentación y evaluación}\par
\vspace{7mm}

\begin{tabularx}{\textwidth}{Y Y Y}
\begin{softcard}[height=33mm,valign=top]{AIEPRed}
\labeltext{Fecha}
{\Large\bfseries\color{AIEPNavy}5 de octubre}\par
Lunes · evento final.
\end{softcard}
&
\begin{softcard}[height=33mm,valign=top]{AIEPCyan}
\labeltext{Participación}
{\Large\bfseries\color{AIEPNavy}10 equipos}\par
6 estudiantes por equipo.
\end{softcard}
&
\begin{softcard}[height=33mm,valign=top]{AIEPGreen}
\labeltext{Evaluación}
{\Large\bfseries\color{AIEPNavy}3 jurados}\par
Tres perspectivas complementarias.
\end{softcard}
\end{tabularx}

\vspace{4mm}
\begin{card}
\labeltext{Datos del evento}
\begin{tabularx}{\textwidth}{@{}L{39mm}Y@{}}
\textbf{Socio comunitario} & Instituto Comercial Liceo Bicentenario, Osorno.\\
\textbf{Instancias previas} & Tres talleres y cuatro mentorías.\\
\textbf{Responsabilidad} & Equipo conjunto AIEP y socio comunitario.\\
\textbf{Resultado} & Propuestas presentadas, evaluación del jurado, retroalimentación, reconocimiento y registro institucional.\\
\end{tabularx}
\end{card}

\begin{softcard}{AIEPBlue}
\labeltext{Propósito}
El concurso constituye el cierre público, pedagógico y competitivo del programa. Cada equipo
presenta una propuesta de innovación sostenible desarrollada a partir de un problema real,
explica las decisiones tomadas, demuestra la evidencia reunida y comunica qué valor podría
aportar su solución.
\end{softcard}

\begin{card}
\labeltext{Trayectoria que se evalúa}
\centering
{\large\bfseries\color{AIEPNavy}
Problema → contexto y personas → solución → tecnología y datos → pruebas y mejoras → comunicación pública}
\end{card}

\begin{softcard}{AIEPGold}
\labeltext{GeoGreen como referente}
GeoGreen demuestra hasta dónde puede crecer una idea, pero no constituye una solución que deba
copiarse. Una propuesta puede ser más simple, diferente o superar el referente, siempre que sus
decisiones sean coherentes y su avance pueda explicarse mediante evidencia.
\end{softcard}

\begin{softcard}{AIEPRed}
\labeltext{Alcance de estos lineamientos}
Este marco regula participación, presentación, jurado, evaluación, desempates, seguridad y
registros. La cantidad de reconocimientos, las categorías y los premios se determinan mediante
la coordinación institucional correspondiente.
\end{softcard}

\newpage
% 2 ---------------------------------------------------------------------------
\kicker{Criterios comunes}
{\LARGE\bfseries\color{AIEPNavy} Principios y condiciones de participación}\par
\vspace{4mm}

\begin{tabularx}{\textwidth}{Y Y}
\begin{card}[height=84mm,valign=top]
\labeltext{Principios del concurso}
\begin{enumerate}
  \item Se evalúa el proceso, no solamente el objeto final.
  \item La evidencia vale más que la apariencia.
  \item La tecnología debe cumplir un propósito.
  \item Las mentorías orientan; el desarrollo pertenece al equipo.
  \item Todos conocen anticipadamente criterios, tiempos y condiciones.
  \item La seguridad está por encima de la demostración.
  \item La comunicación distingue entre lo construido, probado, simulado y proyectado.
  \item La inteligencia artificial puede apoyar; el equipo debe verificar, comprender y defender sus decisiones.
\end{enumerate}
\end{card}
&
\begin{card}[height=84mm,valign=top]
\labeltext{Condiciones comunes}
\begin{itemize}
  \item Participan los diez equipos constituidos durante el programa.
  \item Cada equipo presenta una sola propuesta.
  \item La propuesta se relaciona con sostenibilidad, medioambiente, residuos, reciclaje, uso responsable de recursos o una necesidad comunitaria asociada.
  \item No es obligatorio llegar con un producto industrial terminado.
  \item La madurez alcanzada se valora cuando está respaldada por pruebas, aprendizajes y decisiones verificables.
  \item El costo o la cantidad de componentes no entrega puntaje por sí mismo.
\end{itemize}
\end{card}
\end{tabularx}

\begin{softcard}{AIEPCyan}
\labeltext{Formas válidas de evidencia}
Prototipo físico · simulación · interfaz · maqueta · diagrama · video · fotografía · registro de
prueba · análisis de datos · comparación de versiones · explicación documentada de un error.
\end{softcard}

\begin{card}
\labeltext{Requisitos mínimos de la propuesta}
El equipo debe poder explicar:
\begin{enumerate}
  \item qué problema aborda y en qué contexto ocurre;
  \item a quién afecta o beneficia;
  \item qué solución propone;
  \item qué variable observa y qué función cumple la tecnología;
  \item qué ocurre desde la entrada hasta la respuesta del sistema;
  \item qué construyó, simuló, investigó o probó;
  \item qué resultado obtuvo, qué error encontró o qué decisión mejoró;
  \item qué alcance real tiene actualmente la propuesta.
\end{enumerate}
\end{card}

\begin{softcard}{AIEPGreen}
\labeltext{Participación del equipo}
Las seis personas realizan un aporte sustantivo. La contribución puede expresarse mediante
explicación oral, operación de la demostración, presentación de evidencia, soporte visual,
respuesta a preguntas o coordinación de una parte del relato.
\end{softcard}

\newpage
% 3 ---------------------------------------------------------------------------
\kicker{Presentación pública}
{\LARGE\bfseries\color{AIEPNavy} Formato, soportes y composición del jurado}\par
\vspace{4mm}

\begin{softcard}{AIEPRed}
\labeltext{Bloque máximo por equipo · 8 minutos}
\begin{tabularx}{\textwidth}{@{}L{50mm}C{34mm}Y@{}}
\textbf{Pitch y demostración} & \textbf{4 minutos} & Problema, solución, funcionamiento, evidencia y aporte.\\
\textbf{Preguntas del jurado} & \textbf{2 minutos} & Decisiones, límites, pruebas y aprendizajes.\\
\textbf{Registro y transición} & \textbf{2 minutos} & Cierre de puntajes y preparación del equipo siguiente.\\
\end{tabularx}
\end{softcard}

\begin{card}
\labeltext{Secuencia recomendada del pitch}
\begin{enumerate}
  \item Problema y contexto.
  \item Personas o comunidad relacionada.
  \item Solución propuesta.
  \item Funcionamiento: entrada, procesamiento, comunicación y respuesta.
  \item Evidencia de prueba, error, aprendizaje o mejora.
  \item Aporte esperado y estado actual del proyecto.
\end{enumerate}
\end{card}

\begin{softcard}{AIEPCyan}
\labeltext{Uso de soportes}
El equipo puede utilizar presentación, afiche, imágenes, video breve, diagrama, simulación,
interfaz, maqueta o prototipo. Cada soporte debe ayudar a situar, explicar o demostrar. La
demostración y los videos forman parte de los cuatro minutos del pitch.
\end{softcard}

\begin{tabularx}{\textwidth}{Y Y Y}
\begin{card}[height=58mm,valign=top]
\labeltext{Perfil técnico}
Programación, electrónica, tecnología o innovación.\par
\vspace{3mm}
Observa funcionamiento, datos, evidencia técnica, seguridad y alcance.
\end{card}
&
\begin{card}[height=58mm,valign=top]
\labeltext{Perfil social y ambiental}
Pertinencia, sostenibilidad e impacto comunitario.\par
\vspace{3mm}
Observa problema, personas relacionadas, relevancia y uso responsable de recursos.
\end{card}
&
\begin{card}[height=58mm,valign=top]
\labeltext{Perfil educativo}
Representación educativa o del socio comunitario.\par
\vspace{3mm}
Observa contexto escolar, comprensión, participación y valor formativo.
\end{card}
\end{tabularx}

\begin{card}
\labeltext{Moderación y secretaría · sin derecho a voto}
Una cuarta persona presenta a los equipos, controla los tiempos, ordena las preguntas, resguarda
las pautas, consolida puntajes, aplica el desempate y registra el resultado en el acta.
\end{card}

\begin{softcard}{AIEPGold}
\labeltext{Actuación del jurado}
Cada integrante evalúa individualmente con la misma rúbrica, completa su pauta antes de deliberar,
formula preguntas breves y entrega una fortaleza y una oportunidad de mejora por equipo. Antes de
comenzar se realiza una calibración común de criterios y niveles.
\end{softcard}

\newpage
% 4 ---------------------------------------------------------------------------
\kicker{Evaluación común}
{\LARGE\bfseries\color{AIEPNavy} Rúbrica oficial · 100 puntos}\par
\vspace{4mm}

\begin{tabularx}{\textwidth}{@{}L{55mm}C{18mm}Y@{}}
\rowcolor{AIEPNavy}
\textcolor{White}{\bfseries Criterio} &
\textcolor{White}{\bfseries Peso} &
\textcolor{White}{\bfseries Qué observa}\\
Problema y pertinencia ambiental o comunitaria & 15 & Claridad, contexto, personas relacionadas y relevancia sostenible.\\
\rowcolor{AIEPSoftBlue}
Coherencia entre problema, solución y personas usuarias & 20 & Relación lógica, utilidad, adaptación propia y comprensión de la necesidad.\\
Uso pertinente de tecnología y datos & 15 & Función del sensor, flujo entrada–proceso–salida y propósito técnico.\\
\rowcolor{AIEPGreenSoft}
Evidencia, pruebas, aprendizaje e iteración & 25 & Respaldo, pruebas, resultados, errores reconocidos y mejoras aplicadas.\\
Viabilidad, seguridad y uso responsable de recursos & 10 & Alcance realista, límites, materiales, riesgos y próximos pasos.\\
\rowcolor{AIEPGoldSoft}
Trabajo del equipo y comunicación & 15 & Coordinación, participación, claridad, soporte y respuesta a preguntas.\\
\end{tabularx}

\vspace{5mm}
\begin{card}
\labeltext{Niveles de desempeño}
\begin{tabularx}{\textwidth}{@{}C{14mm}L{31mm}Y@{}}
\level{AIEPGreen}{4} & \textbf{Sólido} & La dimensión es clara, coherente, verificable y está bien integrada.\\
\level{AIEPBlue}{3} & \textbf{Logrado} & Se comprende y cuenta con respaldo suficiente, aunque admite mejoras.\\
\level{AIEPGold}{2} & \textbf{En desarrollo} & Existe una base pertinente, con vacíos o evidencia parcial.\\
\level{AIEPRed}{1} & \textbf{Inicial} & Aparece de forma general y todavía no sostiene la propuesta.\\
\level{AIEPSlate}{0} & \textbf{No observable} & No fue presentada o no existe información para evaluarla.\\
\end{tabularx}
\end{card}

\begin{softcard}{AIEPCyan}
\labeltext{Cálculo}
\centering
{\large\bfseries\color{AIEPNavy}
Puntaje del criterio = (nivel obtenido ÷ 4) × peso del criterio}
\par\vspace{2mm}
El resultado del equipo corresponde al promedio de los puntajes totales asignados por las tres
personas del jurado.
\end{softcard}

\begin{softcard}{AIEPGreen}
\labeltext{Criterio central}
La evidencia, las pruebas y las mejoras concentran el mayor peso. Un proyecto no obtiene ventaja
por verse más complejo: obtiene ventaja cuando demuestra qué hizo, qué aprendió y cómo esa
evidencia mejoró la propuesta.
\end{softcard}

\begin{softcard}{AIEPRed}
\labeltext{Imparcialidad}
No se otorgan ventajas por afinidad, estética, costo, cantidad de componentes o complejidad
aparente. El jurado valora aquello que el equipo presenta, explica y demuestra.
\end{softcard}

\newpage
% 5 ---------------------------------------------------------------------------
\kicker{Descriptores de evaluación}
{\LARGE\bfseries\color{AIEPNavy} Problema, coherencia y tecnología}\par
\vspace{4mm}

\begin{softcard}{AIEPRed}
\labeltext{1 · Problema y pertinencia ambiental o comunitaria · 15 puntos}
\begin{tabularx}{\textwidth}{@{}C{11mm}Y@{}}
\textbf{4} & Problema concreto, situado y relevante; identifica personas o contexto y explica con claridad su relación con sostenibilidad.\\
\textbf{3} & El problema y su relevancia se comprenden; el contexto o sus respaldos admiten mayor profundidad.\\
\textbf{2} & Problema general y pertinente, con contexto, destinatarios o relevancia poco precisos.\\
\textbf{1} & Tema amplio sin delimitación clara de la situación que se busca abordar.\\
\textbf{0} & No presenta un problema reconocible.\\
\end{tabularx}
\end{softcard}

\begin{softcard}{AIEPBlue}
\labeltext{2 · Coherencia entre problema, solución y personas usuarias · 20 puntos}
\begin{tabularx}{\textwidth}{@{}C{11mm}Y@{}}
\textbf{4} & La solución responde directamente al problema, considera a quienes se relacionan con él y muestra decisiones propias justificadas.\\
\textbf{3} & Existe una relación clara, con aspectos de uso o adaptación todavía mejorables.\\
\textbf{2} & La solución se relaciona parcialmente con el problema y mantiene supuestos no resueltos.\\
\textbf{1} & La propuesta se describe, pero no demuestra por qué respondería a la necesidad.\\
\textbf{0} & Problema y solución no guardan una relación observable.\\
\end{tabularx}
\end{softcard}

\begin{softcard}{AIEPCyan}
\labeltext{3 · Uso pertinente de tecnología y datos · 15 puntos}
\begin{tabularx}{\textwidth}{@{}C{11mm}Y@{}}
\textbf{4} & Explica correctamente qué observa, cómo procesa el dato y qué respuesta produce; cada recurso técnico cumple una función necesaria.\\
\textbf{3} & El funcionamiento principal es comprensible y mayormente coherente, con detalles técnicos incompletos.\\
\textbf{2} & Identifica componentes o funciones, pero el flujo del sistema presenta vacíos.\\
\textbf{1} & Enumera tecnología sin explicar cómo contribuye a resolver el problema.\\
\textbf{0} & No presenta uso de tecnología o datos.\\
\end{tabularx}
\end{softcard}

\begin{card}
\labeltext{Pregunta transversal del jurado}
{\large\bfseries\color{AIEPNavy}
¿Se puede seguir el recorrido entre lo que el equipo observó, lo que decidió construir y la
respuesta que espera producir?}
\end{card}

\newpage
% 6 ---------------------------------------------------------------------------
\kicker{Descriptores de evaluación}
{\LARGE\bfseries\color{AIEPNavy} Evidencia, viabilidad y comunicación}\par
\vspace{4mm}

\begin{softcard}{AIEPGreen}
\labeltext{4 · Evidencia, pruebas, aprendizaje e iteración · 25 puntos}
\begin{tabularx}{\textwidth}{@{}C{11mm}Y@{}}
\textbf{4} & Presenta evidencia clara de pruebas, interpreta resultados, reconoce errores y demuestra mejoras derivadas de lo aprendido.\\
\textbf{3} & Muestra al menos una prueba pertinente y explica qué resultado o aprendizaje produjo.\\
\textbf{2} & Presenta evidencia parcial o poco interpretada; su relación con las decisiones es débil.\\
\textbf{1} & Formula afirmaciones sobre el proyecto sin respaldo suficiente.\\
\textbf{0} & No presenta pruebas, registros ni evidencia verificable.\\
\end{tabularx}
\end{softcard}

\begin{softcard}{AIEPGold}
\labeltext{5 · Viabilidad, seguridad y uso responsable de recursos · 10 puntos}
\begin{tabularx}{\textwidth}{@{}C{11mm}Y@{}}
\textbf{4} & Reconoce alcance, límites, recursos y riesgos; la propuesta es segura y plantea una ruta realista.\\
\textbf{3} & La propuesta es realizable y segura, aunque algunos recursos, límites o pasos requieren precisión.\\
\textbf{2} & Mantiene dudas importantes de alcance, recursos o seguridad.\\
\textbf{1} & Depende de supuestos o recursos no explicados y no aborda adecuadamente los riesgos.\\
\textbf{0} & La demostración es insegura o no existe información para analizar la viabilidad.\\
\end{tabularx}
\end{softcard}

\begin{softcard}{AIEPBlue}
\labeltext{6 · Trabajo del equipo y comunicación · 15 puntos}
\begin{tabularx}{\textwidth}{@{}C{11mm}Y@{}}
\textbf{4} & El equipo se coordina, distribuye aportes sustantivos, utiliza el soporte con propósito y responde con claridad y honestidad.\\
\textbf{3} & La presentación es comprensible y coordinada, con oportunidades menores de síntesis, participación o uso del soporte.\\
\textbf{2} & Se comprende la idea principal, pero existen desequilibrios de participación, lectura excesiva o una secuencia poco integrada.\\
\textbf{1} & La presentación dificulta comprender el proyecto o depende de una sola persona.\\
\textbf{0} & El equipo no logra presentar una versión evaluable.\\
\end{tabularx}
\end{softcard}

\begin{card}
\labeltext{Pregunta transversal del jurado}
{\large\bfseries\color{AIEPNavy}
¿El equipo puede demostrar qué hizo, explicar qué aprendió y reconocer con honestidad el estado
actual de su propuesta?}
\end{card}

\newpage
% 7 ---------------------------------------------------------------------------
\kicker{Aplicación de la evaluación}
{\LARGE\bfseries\color{AIEPNavy} Procedimiento, desempate y seguridad}\par
\vspace{4mm}

\begin{tabularx}{\textwidth}{Y Y}
\begin{card}[height=105mm,valign=top]
\labeltext{Procedimiento}
\begin{enumerate}
  \item La moderación identifica al equipo y activa el tiempo.
  \item El equipo realiza su pitch y demostración.
  \item El jurado formula preguntas.
  \item Cada integrante completa su pauta individual sin deliberación.
  \item La secretaría recibe las pautas y consolida los puntajes.
  \item Una vez presentados todos los equipos, se revisan los resultados agregados.
  \item Si existe empate relevante para el reconocimiento, se aplica el procedimiento establecido.
  \item El acta registra puntajes, desempate y resultado final.
\end{enumerate}
\end{card}
&
\begin{card}[height=105mm,valign=top]
\labeltext{Desempate}
Los empates se resuelven en este orden:
\begin{enumerate}
  \item mayor promedio en evidencia, pruebas, aprendizaje e iteración;
  \item mayor promedio en coherencia entre problema, solución y personas usuarias;
  \item mayor promedio en problema y pertinencia ambiental o comunitaria;
  \item deliberación fundada y votación de las tres personas del jurado.
\end{enumerate}
\vspace{3mm}
La aplicación del desempate queda registrada en el acta.
\end{card}
\end{tabularx}

\begin{softcard}{AIEPRed}
\labeltext{Seguridad durante las demostraciones}
\begin{itemize}
  \item La moderación o el perfil técnico puede detener una demostración que implique riesgo.
  \item Una conexión insegura no se energiza durante el evento.
  \item Detener una demostración no elimina la participación: se evalúa la evidencia segura disponible.
  \item No se improvisan cambios eléctricos durante la presentación.
\end{itemize}
\end{softcard}

\begin{softcard}{AIEPCyan}
\labeltext{Contingencias}
Los equipos deben contar con fotografías, video o registros que permitan explicar una prueba si
el prototipo falla. Un fallo en vivo no constituye automáticamente un fracaso: se evalúa la
capacidad de explicar qué ocurrió, qué evidencia previa existe y cómo se enfrentaría el problema.
\end{softcard}

\begin{softcard}{AIEPGold}
\labeltext{Orden y equidad}
Todos los equipos disponen del mismo tiempo máximo, reciben preguntas dentro del mismo bloque y
son evaluados con idénticos criterios. El orden de presentación se comunica antes del inicio.
\end{softcard}

\newpage
% 8 ---------------------------------------------------------------------------
\kicker{Cierre y trazabilidad}
{\LARGE\bfseries\color{AIEPNavy} Retroalimentación, registros y reconocimientos}\par
\vspace{4mm}

\begin{tabularx}{\textwidth}{Y Y}
\begin{card}[height=63mm,valign=top]
\labeltext{Cada equipo recibe}
\begin{itemize}
  \item puntaje consolidado por criterio;
  \item una fortaleza observada;
  \item una oportunidad concreta de mejora.
\end{itemize}
\vspace{3mm}
La devolución describe acciones y evidencias; no califica rasgos personales.
\end{card}
&
\begin{card}[height=63mm,valign=top]
\labeltext{El evento conserva}
\begin{itemize}
  \item pautas individuales del jurado;
  \item acta de consolidación y resultado;
  \item registro de desempate, si se utiliza;
  \item identificación de las propuestas;
  \item fotografías y registro audiovisual institucional.
\end{itemize}
\end{card}
\end{tabularx}

\begin{softcard}{AIEPGreen}
\labeltext{Instrumentos operativos asociados}
\begin{enumerate}
  \item pauta individual de evaluación;
  \item planilla o acta de consolidación de puntajes;
  \item guía breve de jurado y calibración;
  \item ficha de identificación y orden de presentación;
  \item programa operativo del evento;
  \item formato de retroalimentación para los equipos.
\end{enumerate}
\end{softcard}

\begin{softcard}{AIEPRed}
\labeltext{Modalidad de reconocimiento}
La dirección, coordinación y administración responsables determinan la cantidad de equipos
reconocidos, la existencia de posiciones o categorías y las características de los premios. La
modalidad seleccionada se comunica antes del evento y se aplica de manera consistente a los
resultados obtenidos mediante esta evaluación.
\end{softcard}

\begin{card}
\labeltext{Compatibilidad del sistema}
La rúbrica y el procedimiento se mantienen sin modificaciones si se decide reconocer a un equipo,
establecer un podio o entregar reconocimientos por categorías. La evaluación produce un resultado
comparable y trazable; la modalidad de premiación define cómo se traduce ese resultado en
reconocimientos.
\end{card}

\begin{softcard}{AIEPBlue}
\labeltext{Sentido del cierre}
La competencia reconoce una trayectoria de observación, colaboración, desarrollo, prueba y
comunicación. El resultado final permite celebrar los avances y conservar evidencia útil para
fortalecer futuras versiones de GeoGreen Escolar.
\end{softcard}

\begin{tcolorbox}[colback=AIEPNavy,colframe=AIEPNavy,arc=1.2mm,halign=center,top=3mm,bottom=3mm]
{\large\bfseries\color{White}Una propuesta se fortalece cuando puede explicarse, probarse y mejorarse.}
\end{tcolorbox}

\end{document}
